% 以下为 5.7 节局部替换候选，具体插入位置见注释；不包含其他章节修改。

% 替换第15行中从“沿束流轴的轴对称性”至“若存在非对称极化……”的句子：
沿束流轴的轴对称性使统计张量仅有 $q=0$ 分量。在所选量子化基中，若磁子能级布居满足 $P(M_i)=P(-M_i)$，则奇阶取向张量为零；若该对称性不成立，奇阶张量可以非零。对具有确定宇称的核态之间的通常电磁跃迁，未分析光子偏振时需对两个螺旋度求和，强度分布中的奇阶贡献因此消失，即使母态本身具有奇阶取向张量也是如此。对于最高多极阶数 $L_{\max}=2$ 的情形，强度展开最多保留 $k=0,2,4$，并还受 $k\leq 2I_i$ 限制。

% 第13行沿用现有定义；如讨论奇阶项，紧接该定义增加：
式中 CG 系数的次序规定了张量相位；与 Hamilton 式~(12.73) 的张量记号比较时，按此定义的奇阶分量相差一个负号，偶阶分量相同。以下未分析偏振的强度式只使用偶阶分量。

% 替换第16--18行之后的一句说明：
若初态经一条指定的、未被探测的级联达到 $I_i$，并采用所述退取向传递近似，$B_k(I_i)=B_k(J_{\mathrm{CN}})\prod_m U_k(m)$；对多条可能路径，应以相应布居分支和传递系数加权。后文若已使用目标态的 $B_k(I_i)$，则不再重复乘入相同的 $U_k$。

% 替换第36行：
其中表观展开系数为
\[
a_k=\frac{B_k(I_i)Q_k A_k(\delta)}{B_0(I_i)Q_0 A_0}
    =B_k(I_i)Q_k A_k(\delta),\qquad k=2,4,
\]
第二个等号采用 $B_0=Q_0=A_0=1$ 的归一约定。$Q_k$ 是以非负探测效率加权、按零阶响应归一的轴对称角接受因子；若实际响应不满足这一简化，应直接积分角响应。

% 替换第42行中“采用 Rose-Brink 规范体系”为：
采用 Krane--Steffen（KS）混合比与 Hamilton 所列 $F_k$ 的相位约定，并与约化矩阵元所用的算符单位及相位保持一致~\cite{lange-kumar-hamilton-1982-multipole-admixtures,book-hamilton-electromagnetic-interaction}。
% 第44行的 sqrt(3)/10 关系沿用，但应确认全文磁算符实际采用的单位约定；若无法确认，不以该式宣称已完成数值/符号换算。
% 在第54行后补充：
在 Lange 等给出的发射约定下，$\delta_{\rm KS}=-\delta_{\rm RB}$；具体转换仍须连同相应几何系数、初末态排列及算符定义使用。式~\eqref{eq:5-7-ak-rational} 的正干涉项对应所采用的 KS 约定，不应把它直接标成 RB 公式。

% 第72式后可补充：
式~\eqref{eq:5-7-dco-def} 使用同一套计数、门宽、曝光和本底扣除约定，并假设角组的双探测效率可由所列单探测效率乘积表示；若角组由多对探测器构成或存在相关响应，需以实际探测器对的响应加权。门跃迁可位于待测跃迁之前或之后，具体关联应按级联时间次序计算。

% 替换第106--111行的偏振理论解释（原式标签保留）：
在所约定的线偏振基底下，令 $W_\parallel$ 与 $W_\perp$ 为沿反应平面和平面法向的正交电矢量分量的强度，则
\begin{equation}\label{eq:5-7-pol-theor}
P(\theta)=\frac{W_\parallel-W_\perp}{W_\parallel+W_\perp}
=\frac{\sum_{k\geq2,\,\mathrm{even}}B_k A_k^{\mathrm{pol}}(\delta)Q_k^{\mathrm{pol}}P_k^{2}(\cos\theta)}{\sum_{k\,\mathrm{even}}B_k A_k(\delta)Q_kP_k(\cos\theta)}.
\end{equation}
这里 $A_k^{\mathrm{pol}}$ 包含所选伴随 Legendre 函数的归一、偏振基底及电磁相位；若采用 Hamilton 的 $A_{k2}$，还须包含 $2\sqrt{(k-2)!/(k+2)!}$。$Q_k^{\mathrm{pol}}$ 表示与强度分布相容的有限接受修正；具体响应应结合探测器模型确定。式~\eqref{eq:5-7-asym-exp} 的散射方向、校正系数和灵敏度 $Q(E_\gamma)$ 应使用同一标定的正负号；后者是偏振计灵敏度，与 $Q_k$ 不是同一量。偏振度只有在分母非零且灵敏度可确定时才能提取，混合跃迁的正偏振本身并不证明其为纯 $E2$。

% 替换第123行并修订图注中的“必须”：
式~\eqref{eq:5-7-quadratic-roots} 应连同原比值的非零分母共同求解。二次系数为零时可退化为一次方程；若全部系数为零，则原比值定义域内每个有限 $\delta$ 均满足该等式；若仅剩非零常数则无有限根。令 $\delta=u/v$，纯高阶多极是 $v=0$ 的端点；当 $d_2\ne0$ 且 $R_{\rm DCO}=n_2/d_2$ 时，应另行计入该端点。它不同于有限 $\delta$ 处的分母零点，后者必须检查原始比值定义域。两个端点 $\arctan\delta=\pm90^\circ$ 代表同一纯高阶振幅方向。偏振或内转换测量只有在不确定度允许区分各分支时才能解除简并。

% 第134行图注建议：将“测量误差范围”改为“假设观测带”；将“双螺旋度正定光子发射响应”改为“满足 W±H>0 的示意二次函数”；明确 DCO 与偏振系数均为人为选取，未指定核态或仪器；“必须引入……方可”改为“可借助具有足够区分力的独立观测仲裁”。producer 图例 Measured 也改成 Hypothetical。
